\section{Conclusion}

We investigated whether small language models can learn matrix inversion through supervised LoRA fine-tuning on synthetic data---a diagnostic case study, not a comprehensive algorithmic benchmark. Our three-level evaluation framework (syntactic, structural, functional) reveals a nuanced outcome.

\textbf{What succeeded}: Among parseable fine-tuned outputs, several coarse numeric diagnostics improved relative to the baseline. The determinant ratio rose from 0.045 to 0.639, sign agreement from 0.503 to 0.795, and relative Frobenius error to the gold inverse fell from 1.030 to 0.727. Parseable fine-tuned outputs were also substantially farther from trivial input copying ($\|X - A\|_F$: 1.45 $\to$ 3.18). These shifts indicate partial inverse-like structure in the fine-tuned model's parseable outputs, going beyond mere format imitation.

\textbf{What failed}: True inverse accuracy remained 0\% under all residual-based tolerances for both models. Fine-tuning also introduced severe generation instability: only 60\% of fine-tuned outputs were parseable (vs.\ 100\% baseline), with exactly half of $3 \times 3$ examples failing, distributed between repetition loops (22\%) and malformed shapes (18\%).

\textbf{Dominant bottlenecks}: The experiment reveals two compounding issues. First, \emph{insufficient numeric precision}: the model shows improved coarse alignment but cannot satisfy the algebraic constraint $AX \approx I$ at any tolerance. Second, \emph{generation instability}: every fine-tuned output reached the 4{,}096-token inference ceiling, indicating a breakdown in length-controlled termination that prevents reliable evaluation of numeric quality. Addressing reliability is a prerequisite before further numeric progress can be meaningfully assessed.

\textbf{Caveats}: All structural improvements are measured conditionally over the 60\% of fine-tuned outputs that are parseable; direct comparison with the fully parseable baseline is not strictly equivalent. Results represent a single training run without variance analysis.

\textbf{Outlook}: This result is a diagnostic, not a verdict. Future work should prioritize: (1) simplifying the output serialization format to reduce format-induced failures; (2) decoding stabilization to address looping; (3) small-scale overfitting experiments to separate model capacity limits from format and reliability bottlenecks; and (4) intermediate-step supervision that provides signal about the computation process, not only the final answer. Scaling experiments, model comparisons, and hyperparameter search are best pursued after reliability is established.
